% Reference-list transcription from the original PDF.
% NOT a verified bibliography database.
% Apparent extraction errors are retained and flagged for checking.

\section*{References}

\begingroup
\small
\setlength{\parindent}{0pt}
\setlength{\parskip}{0.6em}

\hangindent=1.5em\hangafter=1
Barrett, D., Hill, F., Santoro, A., Morcos, A., \& Lillicrap, T.
(2018, July).
Measuring abstract reasoning in neural networks.
In International conference on machine learning (pp. 511-520).
PMLR.
\par

\hangindent=1.5em\hangafter=1
Battaglia, P. W., Hamrick, J. B., \& Tenenbaum, J. B. (2013).
Simulation as an engine of physical scene understanding.
Proceedings of the National Academy of Sciences, 110(45), 18327-18332.
\par

\hangindent=1.5em\hangafter=1
Blalock, D., Ortiz, J. J. G., Frankle, J., \& Guttag, J. (2020).
What is the state of neural network pruning?
arXiv preprint arXiv:2003.03033.
\par

% CHECK: probable title extraction error ("Checs").
\hangindent=1.5em\hangafter=1
Blanch, A. (2020).
Checs and individual differences.
Cambridge University Press.
\par

\hangindent=1.5em\hangafter=1
Bongard, M. (1967).
The problems of knowledge formation.
\par

\hangindent=1.5em\hangafter=1
Botvinick, M., Ritter, S., Wang, J. X., Kurth-Nelson, Z.,
Blundell, C., \& Hassabis, D. (2019).
Reinforcement learning, fast and slow.
Trends in Cognitive Sciences, 23(5), 408-422.
\par

\hangindent=1.5em\hangafter=1
Botvinick, M., Wang, J. X., Dabney, W., Miller, K. J.,
\& Kurth-Nelson, Z. (2020).
Deep reinforcement learning and its neuroscientific implications.
Neuron, 107(4), 603-616.
\par

\hangindent=1.5em\hangafter=1
Bringsjord, S., \& Schimanski, B. (2003, August).
What is artificial intelligence? Psychometric AI as an answer.
In International joint conference on artificial intelligence
(pp. 887-893).
\par

% CHECK: incomplete author list and extraction punctuation.
\hangindent=1.5em\hangafter=1
Brown, T. B., Mann, B., Ryder, N., Subbiah, M., Kaplan, J.,
Dhariwal, P., . Amodei, D. (2020).
Language models are few-shot learners.
arXiv preprint arXiv:2005.14165.
\par

\hangindent=1.5em\hangafter=1
Burda, Y., Edwards, H., Storkey, A., \& Klimov, O. (2018).
Exploration by random network distillation.
arXiv preprint arXiv:1810.12894.
\par

\hangindent=1.5em\hangafter=1
Burgoyne, A. P., \& Engle, R. W. (2020).
Attention control: A cornerstone of higher-order cognition.
Current Directions in Psychological Science, 29(6), 624-630.
\par

\hangindent=1.5em\hangafter=1
Carpenter, P. A., Just, M. A., \& Shell, P. (1990).
What one intelligence test measures: A theoretical account of the
processing in the Raven Progressive Matrices Test.
Psychological Review, 97(3), 404.
\par

% CHECK: author-list completeness.
\hangindent=1.5em\hangafter=1
Ching, T., Himmelstein, D. S., Beaulieu-Jones, B. K.,
Kalinin, A. A., Do, B. T., Way, G. P., Greene, C. S. (2018).
Opportunities and obstacles for deep learning in biology and medicine.
Journal of the Royal Society Interface, 15(141), 20170387.
\par

% CHECK: author extraction; main text cites Chollet.
\hangindent=1.5em\hangafter=1
Choller, F. (2019).
On the measure of intelligence.
arXiv preprint arXiv:1911.01547.
\par

\hangindent=1.5em\hangafter=1
Cole, D. (2020).
The Chinese room argument.
In E. N. Zalta (Ed.),
The Stanford encyclopedia of philosophy (Winter 2020 ed.).
\par

\hangindent=1.5em\hangafter=1
Correa, W. F., Prade, H., \& Richard, G. (2012, August).
When intelligence is just a matter of copying.
In, vol. 12.
European conference on artificial intelligence (pp. 276-281).
\par

\hangindent=1.5em\hangafter=1
Cronbach, L. J. (1957).
The two disciplines of scientific psychology.
American Psychologist, 12(11), 671.
\par

\hangindent=1.5em\hangafter=1
de Groot, A. D. (1946).
Het denken van den schaker: een experimenteel psychologische studie
(p. 315).
Amsterdam: Noord-Hollandische Uitgevers Maatschappij.
\par

\hangindent=1.5em\hangafter=1
Deng, J., Dong, W., Socher, R., Li, L. J., Li, K.,
\& Fei-Fei, L. (2009, June).
Imagenet: A large-scale hierarchical image database.
In 2009 IEEE conference on computer vision and pattern recognition
(pp. 248-255).
IEEE.
\par

\hangindent=1.5em\hangafter=1
Dennett, D. C. (2006).
The frame problem of AI.
Philosophy of Psychology: Contemporary Readings, 433, 67-83.
\par

\hangindent=1.5em\hangafter=1
Dreyfus, H. (1965).
Alchemy and AI.
Santa Monica, CA: RAND Corporation.
\par

\hangindent=1.5em\hangafter=1
Dubhashi, D., \& Lappin, S. (2017).
AI dangers: Imagined and real.
Communications of the ACM, 60(2), 43-45.
\par

\hangindent=1.5em\hangafter=1
Evans, T. G. (1964).
A program for the solution of a class of geometric analogy
intelligence-test questions.
Bedford, MS: Air Force Cambridge Research Labs.
\par

% CHECK: dissertation title and URL extraction.
% Main text names the architecture Phaeaco.
\hangindent=1.5em\hangafter=1
Foundalis, H. E. (2006).
Pleasure: A cognitive architecture inspired by bongard's problems.
Doctoral dissertation.
Indiana University.
\url{https://www.foundalis.com/res/diss_research}.
\par

\hangindent=1.5em\hangafter=1
Frischkorn, G. T., \& Schubert, A. L. (2018).
Cognitive models in intelligence research: Advantages and
recommendations for their application.
Journal of Intelligence, 6(3), 34.
\par

\hangindent=1.5em\hangafter=1
Gershman, S. J., Pesaran, B., \& Daw, N. D. (2009).
Human reinforcement learning subdivides structured action spaces
by learning effector-specific values.
Journal of Neuroscience, 29(43), 13524-13531.
\par

\hangindent=1.5em\hangafter=1
Gierl, M. J., \& Haladyna, T. M. (2012).
Automatic item generation: Theory and practice.
Routledge.
\par

% CHECK: garbled/incomplete chapter reference.
\hangindent=1.5em\hangafter=1
Gobet, F., \& Charness, N. (2018).
In K. A. Ericsson, R. B. Hoffman, A. Kozbelt,
\& A. M. Williams (Eds.),
Expertise in item: Cambridge handbooks in psychology.
\par

\hangindent=1.5em\hangafter=1
Gobet, F., \& Chassy, P. (2009).
Expertise and intuition: A tale of three theories.
Minds and Machines, 19(2), 151-180.
\par

\hangindent=1.5em\hangafter=1
GPT-3. (2020, September 8).
A robot wrote this entire article. Are you scared yet, human?
The Guardian.
\par

\hangindent=1.5em\hangafter=1
Güçlü, U., \& van Gerven, M. A. (2015).
Deep neural networks reveal a gradient in the complexity of neural
representations across the ventral stream.
Journal of Neuroscience, 35(27), 10005-10014.
\par

\hangindent=1.5em\hangafter=1
He, K., Zhang, X., Ren, S., \& Sun, J. (2016).
Deep residual learning for image recognition.
In Proceedings of the IEEE conference on computer vision and
pattern recognition (pp. 770-778).
\par

\hangindent=1.5em\hangafter=1
Henderson, J. (2020).
The unstoppable rise of computational linguistics in deep learning.
arXiv preprint arXiv:2005.06420.
\par

% CHECK: author extraction; compare Martínez-Plumed and Dowe
% in the main text.
\hangindent=1.5em\hangafter=1
Hernández-Orallo, J., Martínez-Planeel, F., Schmid, U.,
Siebers, M., \& Howe, D. L. (2016).
Computer models solving intelligence test problems:
Progress and implications.
Artificial Intelligence, 230, 74-107.
\par

% CHECK: probable title extraction errors.
\hangindent=1.5em\hangafter=1
Hofstadter, D. (1979).
In Godal, Escher, Buch: an eternal golden braid.
New York, NY: Basic Books.
\par

\hangindent=1.5em\hangafter=1
Ichien, N., Lu, H., \& Holyoak, K. J. (2020).
Verbal analogy problem sets: An inventory of testing materials.
Behavior Research Methods, 52, 1803-1816.
\par

\hangindent=1.5em\hangafter=1
Jobin, A., Ienca, M., \& Vayena, E. (2019).
The global landscape of AI ethics guidelines.
Nature Machine Intelligence, 1(9), 389-399.
\par

% CHECK: initials and spelling; compare names in main text.
\hangindent=1.5em\hangafter=1
Kan, R. J., Ploeger, A., Rajimakers, M. E., Dolan, C. V.,
\& Van Der Maas, H. L. (2010).
Nonlinear epigenetic variance: Review and simulations.
Developmental Science, 13(1), 11-27.
\par

\hangindent=1.5em\hangafter=1
Karras, T., Aila, T., Laine, S., \& Lehtinen, J. (2017).
Progressive growing of gans for improved quality, stability,
and variation.
arXiv preprint arXiv:1710.10196.
\par

\hangindent=1.5em\hangafter=1
Kaul, V., Enslin, S., \& Gross, S. A. (2020).
The history of artificial intelligence in medicine.
Gastrointestinal Endoscopy, 92(4), 807-812.
\par

\hangindent=1.5em\hangafter=1
Khaligh-Razavi, S. M., \& Kriegeskorte, N. (2014).
Deep supervised, but not unsupervised, models may explain
IT cortical representation.
PLoS Computational Biology, 10(11), Article e1003915.
\par

\hangindent=1.5em\hangafter=1
Kiran, B. R., Sobh, I., Talpaert, V., Mannion, P.,
Sallab, A. A. A., Yugamani, S., \& Pérez, P. (2020).
Deep reinforcement learning for autonomous driving: A survey.
arXiv preprint arXiv:2002.00444.
\par

\hangindent=1.5em\hangafter=1
Kriegeskorte, N. (2015).
Deep neural networks: A new framework for modeling biological
vision and brain information processing.
Annual Review of Vision Science, 1, 417-446.
\par

\hangindent=1.5em\hangafter=1
Krizhevsky, A., Sutskever, I., \& Hinton, G. E. (2012).
Imagenet classification with deep convolutional neural networks.
Advances in neural information processing systems, 25, 1097-1105.
\par

\hangindent=1.5em\hangafter=1
Kunda, M., McGreggor, K., \& Goel, A. K. (2013).
A computational model for solving problems from the Raven's
Progressive Matrices intelligence test using iconic visual
representations.
Cognitive Systems Research, 22, 47-66.
\par

\hangindent=1.5em\hangafter=1
Lake, B. M., Ullman, T. D., Tenenbaum, J. B.,
\& Gershman, S. J. (2017).
Building machines that learn and think like people.
Behavioral and Brain Sciences, 40.
\par

\hangindent=1.5em\hangafter=1
Lane, D. M., \& Chang, Y. H. A. (2018).
Chess knowledge predicts chess memory even after controlling
for chess experience: Evidence for the role of high-level processes.
Memory \& Cognition, 46(3), 337-348.
\par

\hangindent=1.5em\hangafter=1
LeCun, Y., Bengio, Y., \& Hinton, G. (2015).
Deep learning.
Nature, 521(7553), 436-444.
\par

\hangindent=1.5em\hangafter=1
Legg, S., \& Hutter, M. (2007).
A collection of definitions of intelligence.
Frontiers in Artificial Intelligence and Applications, 157, 17.
\par

\hangindent=1.5em\hangafter=1
Lindsay, G. W. (2020).
Attention in psychology, neuroscience, and machine learning.
Frontiers in Computational Neuroscience, 14, 29.
\par

\hangindent=1.5em\hangafter=1
Linsley, D., Shiebler, D., Eberhardt, S., \& Serre, T. (2018).
Learning what and where to attend.
arXiv preprint arXiv:1805.08819.
\par

\hangindent=1.5em\hangafter=1
Liu, Y., He, F., Zhang, H., Rao, G., Feng, Z., \& Zhou, Y. (2019).
How well do machines perform on IQ tests:
A comparison study on a large-scale dataset.
In International joint conference on artificial intelligence
(pp. 6110-6116).
\par

\hangindent=1.5em\hangafter=1
Lovett, A., \& Forbus, K. (2012).
Modeling multiple strategies for solving geometric analogy problems.
In, vol. 34.
Proceedings of the annual meeting of the cognitive science society.
No. 34.
\par

\hangindent=1.5em\hangafter=1
Lu, K., Grover, A., Abbeel, P., \& Mordatch, I. (2021).
Pretrained transformers as universal computation engines.
arXiv preprint arXiv:2103.05247.
\par

\hangindent=1.5em\hangafter=1
van der Maas, H. L., Kan, K. J., \& Borsboom, D. (2014).
Intelligence is what the intelligence test measures. Seriously.
Journal of Intelligence, 2(1), 12-15.
\par

% CHECK: author extraction; main text names Marsman.
\hangindent=1.5em\hangafter=1
van der Maas, H. L., Kan, K. J., Borseman, M.,
\& Stevenson, C. E. (2017).
Network models for cognitive development and intelligence.
Journal of Intelligence, 5(2), 16.
\par

\hangindent=1.5em\hangafter=1
van der Maas, H. L., \& Wagenmakers, E. J. (2005).
A psychometric analysis of chess expertise.
The American Journal of Psychology, 29-60.
\par

\hangindent=1.5em\hangafter=1
Mao, J., Gan, C., Kohli, P., Tenenbaum, J. B., \& Wu, J. (2019).
The neuro-symbolic concept learner:
Interpreting scenes, words, and sentences from natural supervision.
arXiv preprint arXiv:1904.12584.
\par

\hangindent=1.5em\hangafter=1
Marblestone, A. H., Wayne, G., \& Kording, K. P. (2016).
Toward an integration of deep learning and neuroscience.
Frontiers in Computational Neuroscience, 10, 94.
\par

% CHECK: author extraction; main text cites Marcus.
\hangindent=1.5em\hangafter=1
Marcos, G. (2018).
Deep learning: A critical appraisal.
arXiv preprint arXiv:1801.00631.
\par

\hangindent=1.5em\hangafter=1
Marr, D. (1982).
Vision: A computational approach.
Freeman \& Co.
\par

% CHECK: journal volume/details.
\hangindent=1.5em\hangafter=1
McCulloch, W. S., \& Pitts, W. (1943).
A logical calculus of the ideas immanent in nervous activity.
The Bulletin of Mathematical Biophysics, 3(4), 115-133.
\par

% CHECK: author-list extraction/completeness.
\hangindent=1.5em\hangafter=1
Mnih, V., Kavukcuoglu, K., Silver, D., Rusu, A. A.,
Veness, J., Bellemare, M. G., Frenkelis, D. (2015).
Human-level control through deep reinforcement learning.
Nature, 518(7540), 529-533.
\par

% CHECK: author-list completeness.
\hangindent=1.5em\hangafter=1
Naselaris, T., Bassett, D. S., Fletcher, A. K., Kording, K.,
Kriegeskorte, N., Nienborg, H., Kay, K. (2018).
Cognitive computational neuroscience:
A new conference for an emerging discipline.
Trends in Cognitive Sciences, 22(5), 365-367.
\par

% CHECK: author extraction; main text cites Neftci.
\hangindent=1.5em\hangafter=1
Nefici, E. O., \& Averbeck, B. B. (2019).
Reinforcement learning in artificial and biological systems.
Nature Machine Intelligence, 1(3), 133-143.
\par

\hangindent=1.5em\hangafter=1
Newell, A. (1973).
You can't play 20 questions with nature and win:
Projective comments on the papers of this symposium.
In W. G. Chase (Ed.), Visual information processing.
New York, NY: Academic Press.
\par

\hangindent=1.5em\hangafter=1
Newell, A., \& Simon, H. (1956).
The logic theory machine---a complex information processing system.
IRE Transactions on Information Theory, 2(3), 61-79.
\par

\hangindent=1.5em\hangafter=1
Newell, A., \& Simon, H. A. (2007).
Computer science as empirical inquiry: Symbols and search.
In ACM turing award lectures (p. 1975).
\par

\hangindent=1.5em\hangafter=1
Nie, W., Yu, Z., Mao, L., Patel, A. R., Zhu, Y.,
\& Anandkumar, A. (2020).
BONGARD-LOGO: A new benchmark for human-level concept learning
and reasoning.
arXiv preprint arXiv:2010.00763.
\par

\hangindent=1.5em\hangafter=1
O'Doherty, J. P., Lee, S. W., \& McNamee, D. (2015).
The structure of reinforcement-learning mechanisms in the human brain.
Current Opinion in Behavioral Sciences, 1, 94-100.
\par

\hangindent=1.5em\hangafter=1
Pearl, J. (2019).
The seven tools of causal inference, with reflections on
machine learning.
Communications of the ACM, 62(3), 54-60.
\par

\hangindent=1.5em\hangafter=1
Penrose, R. (1994).
Shadows of the Mind.
Oxford: Oxford University Press.
\par

\hangindent=1.5em\hangafter=1
Pekar, N., Benny, Y., \& Wolf, L. (2020).
Generating correct answers for progressive matrices intelligence tests.
arXiv preprint arXiv:2011.00496.
\par

% CHECK: initials and chapter/book formatting.
\hangindent=1.5em\hangafter=1
Rescorla, B. A., \& Wagner, A. R. (1972).
In A. H. Black, \& W. F. Prokasy (Eds.),
A theory of Pavlovian conditioning:
Variations in the effectiveness of reinforcement and nonreinforcement,
classical conditioning II (pp. 64-99).
Appleton-Century-Crofts.
\par

\hangindent=1.5em\hangafter=1
Rosenblatt, F. (1958).
The perceptron: A probabilistic model for information storage and
organization in the brain.
Psychological Review, 65(6), 386.
\par

\hangindent=1.5em\hangafter=1
Ruder, S. (2017).
An overview of multi-task learning in deep neural networks.
arXiv preprint arXiv:1706.05098.
\par

\hangindent=1.5em\hangafter=1
Rudin, C. (2019).
Stop explaining black box machine learning models for high stakes
decisions and use interpretable models instead.
Nature Machine Intelligence, 1(5), 206-215.
\par

\hangindent=1.5em\hangafter=1
Schank, R. C. (1980).
How much intelligence is there in artificial intelligence?
Intelligence, 4, 1-14.
\par

\hangindent=1.5em\hangafter=1
Scholte, H. S. (2018).
Fantastic DNimals and where to find them.
Neuroimage, 180, 112-113.
\par

\hangindent=1.5em\hangafter=1
Schrimpf, M., Kubilius, J., Hong, H., Majaj, N. J.,
Rajalingham, R., Issa, E. B., ... DiCarlo, J. J. (2018).
Brain-score: Which artificial neural network for object recognition
is most brain-like?
BioRxiv, 407007.
\par

\hangindent=1.5em\hangafter=1
Schultz, W., Dayan, P., \& Montague, P. R. (1997).
A neural substrate of prediction and reward.
Science, 275(5306), 1593-1599.
\par

% CHECK: publication metadata.
\hangindent=1.5em\hangafter=1
Searle, J. R. (1980).
Minds, brains, and programs.
Verbal Behaviour as the Hallmark of Intelligence.
The Turing Test, 201-224.
\par

\hangindent=1.5em\hangafter=1
Siciliano, B., \& Khatib, O. (Eds.). (2016).
Springer handbook of robotics.
Springer.
\par

% CHECK: author-list extraction.
\hangindent=1.5em\hangafter=1
Silver, D., Hubert, T., Schrittwieser, J., Antonoglou, I.,
Lai, M., Gura, A., ... Hassabis, D. (2018).
A general reinforcement learning algorithm that masters chess,
shogi, and go through self-play.
Science, 362(6419), 1140-1144.
\par

\hangindent=1.5em\hangafter=1
Sutton, R. S., \& Barto, A. G. (1987, July).
A temporal-difference model of classical conditioning.
In Proceedings of the ninth annual conference of the cognitive
science society (pp. 355-378).
\par

% CHECK: book title and numeric detail.
\hangindent=1.5em\hangafter=1
Sutton, R. S., \& Barto, A. G. (1998).
Introduction to reinforcement learning, 135.
Cambridge: MIT press.
\par

\hangindent=1.5em\hangafter=1
Tenenbaum, J. B., Kemp, C., Griffiths, T. L., \& Goodman, N. D.
(2011).
How to grow a mind: Statistics, structure, and abstraction.
Science, 331(6022), 1279-1285.
\par

\hangindent=1.5em\hangafter=1
Torrey, L., \& Shavlik, J. (2010).
Transfer learning.
In E. Soria, J. Martin, R. Magdalena, M. Martinez,
\& A. Serrano (Eds.),
Handbook of research on machine learning applications and trends:
algorithms, methods, and techniques (pp. 242-264).
IGI global.
\par

\hangindent=1.5em\hangafter=1
Van Harreveld, F., Wagenmakers, E. J., \& Van Der Maas, H. L.
(2007).
The effects of time pressure on chess skill:
An investigation into fast and slow processes underlying
expert performance.
Psychological Research, 71(5), 591-597.
\par

\hangindent=1.5em\hangafter=1
VanRullen, R. (2017).
Perception science in the age of deep neural networks.
Frontiers in Psychology, 8, 142.
\par

% CHECK: author-list completeness.
\hangindent=1.5em\hangafter=1
Vaswani, A., Shazeer, N., Parmar, N., Uszkoreit, J.,
Jones, L., Gomez, A. N., Polosukhin, I. (2017).
Attention is all you need.
arXiv preprint arXiv:1706.03762.
\par

\hangindent=1.5em\hangafter=1
Vinyals, O., Blundell, C., Lillicrap, T., Kavukcuoglu, K.,
\& Wierstra, D. (2016).
Matching networks for one shot learning.
arXiv preprint arXiv:1606.04080.
\par

\hangindent=1.5em\hangafter=1
Wang, J. X. (2021).
Meta-learning in natural and artificial intelligence.
Current Opinion in Behavioral Sciences, 38, 90-95.
\par

% CHECK: author-list completeness.
\hangindent=1.5em\hangafter=1
Wang, J. X., Kurth-Nelson, Z., Kumaran, D., Tirumala, D.,
Soyer, H., Leibo, J. Z., Botvinick, M. (2018).
Prefrontal cortex as a meta-reinforcement learning system.
Nature Neuroscience, 21(6), 860-868.
\par

\hangindent=1.5em\hangafter=1
Xu, K., Ba, J., Kiros, R., Cho, K., Courville, A.,
Salakhudinov, R., ... Bengio, Y. (2015, June).
Show, attend and tell: Neural image caption generation
with visual attention.
In International conference on machine learning (pp. 2048-2057).
PMLR.
\par

\hangindent=1.5em\hangafter=1
Yamins, D. L., Hong, H., Cadieu, C. F., Solomon, E. A.,
Seibert, D., \& DiCarlo, J. J. (2014).
Performance-optimized hierarchical models predict neural responses
in higher visual cortex.
Proceedings of the National Academy of Sciences, 111(23), 8619-8624.
\par

\hangindent=1.5em\hangafter=1
Yi, K., Wu, J., Gan, C., Torralba, A., Kohli, P.,
\& Tenenbaum, J. B. (2018).
Neural-symbolic vqa: Disentangling reasoning from vision
and language understanding.
arXiv preprint arXiv:1810.02338.
\par

\hangindent=1.5em\hangafter=1
Yosinski, J., Clune, J., Bengio, Y., \& Lipson, H. (2014).
How transferable are features in deep neural networks?
arXiv preprint arXiv:1411.1792.
\par

\hangindent=1.5em\hangafter=1
Zhang, C., Gao, F., Jia, B., Zhu, Y., \& Zhu, S. C. (2019).
Raven: A dataset for relational and analogical visual reasoning.
In Proceedings of the IEEE/CVF conference on computer vision
and pattern recognition (pp. 5317-5327).
\par

\hangindent=1.5em\hangafter=1
Zhang, L., Zhang, L., \& Du, B. (2016).
Deep learning for remote sensing data:
A technical tutorial on the state of the art.
IEEE Geoscience and Remote Sensing Magazine, 4(2), 22-40.
\par

\hangindent=1.5em\hangafter=1
Zou, J., Huss, M., Abid, A., Mohammadi, P., Torkamani, A.,
\& Telenti, A. (2019).
A primer on deep learning in genomics.
Nature Genetics, 51(1), 12-18.
\par

\endgroup
